\begin{frame}{Lecciones para nuestros proyectos}
\begin{enumerate}
\item \textbf{PresustentacionesAWPS:} un mismo servicio de negocio puede tener dos contratos, REST y SOAP, sin duplicar la lógica. También hay que usar bien los códigos HTTP: cambiamos el \texttt{200} vacío por un \texttt{404}.
\item \textbf{SGIRI:} todo servicio externo necesita un tiempo límite y un plan B.
\item \textbf{Gestión documental:} métodos idempotentes y códigos HTTP correctos permiten reintentar sin duplicar documentos.
\end{enumerate}
\vspace{10pt}
\begin{block}{Conclusión}
SOAP y REST no compiten en todo: el contexto decide. Y ningún servicio web oculta del todo sus fallos, así que el cliente siempre debe estar preparado.
\end{block}
\end{frame}
